\section{Unbinned background subtraction by negative-weight injection}
\label{app:negweight}

The background subtraction of \S\ref{sec:method-algo} is the one step in which
the analysis binning enters before the final histogramming: each data event in
reconstructed bin $b$ is down-weighted by the purity
$w_{\mathrm{pur}}(b)=\max(0,\,N_{\mathrm{data}}(b)-N_{\mathrm{bkg}}(b))/
N_{\mathrm{data}}(b)$. This appendix documents the fully unbinned realization of
the same correction --- injecting the background simulation onto the measured
side of the fit with negative weight --- shows that it targets an identical
estimand, and quantifies its agreement with the binned correction on the 2D
measurement. It is the continuous version of the purity weight recommended for
nontrivial backgrounds by Ref.~\cite{unbinnedguide}, and the natural treatment
for the extended-fiducial analysis, where the background is larger.

\paragraph{Provenance gate on this appendix's numbers.}
\textbf{The twelve real-data/production comparison values in this appendix are STRUCK
(\texttt{OI-130}, gated 2026-08-21); the four synthetic-toy values are ATTESTED and stand.} The
analytic derivation, the estimand identity and the displacement argument need none of them.

\emph{Struck --- twelve real-data/production values} ---
seven central-comparison (\texttt{\textbackslash nwSigPur}, \texttt{\textbackslash nwSigNeg},
\texttt{\textbackslash nwRatioTot}, \texttt{\textbackslash nwPctTot},
\texttt{\textbackslash nwMedianBin}, \texttt{\textbackslash nwRmsBin},
\texttt{\textbackslash nwWorstBin}), two covariance (\texttt{\textbackslash nwSystRatio},
\texttt{\textbackslash nwStatRatio}) and three exact-estimator
(\texttt{\textbackslash nwSpMedian}, \texttt{\textbackslash nwSpWithin},
\texttt{\textbackslash nwSpRawMax}) --- together with the two \emph{derived} residual differences
(\texttt{\textbackslash nwSystResid}, \texttt{\textbackslash nwStatResid}), which are defined from
their two covariance parents so that they cannot outlive a retraction of those parents. All rest on
\emph{untracked scratch products}: \texttt{.gitignore} carries a blanket \texttt{*.root} and
\texttt{git ls-files} returns zero \texttt{.root} files anywhere under \texttt{2d-unfolding/}, so
none of the eight ROOTs in \texttt{HANDOFF\_bkg\_negweight/runs/} nor the 51 bootstrap and 188
universe replicas behind the two covariance ratios is under version control, and no
\texttt{RUNS.tsv} or \texttt{VALIDATION\_LEDGER.md} row cites any of them. Recovering these needs a
\emph{durability} step, not a re-derivation.

\emph{Attested --- four synthetic-toy values} (\texttt{\textbackslash nwToyNeg},
\texttt{\textbackslash nwToyPur}, \texttt{\textbackslash nwToyBias},
\texttt{\textbackslash nwToySeed}) come from the closed-form toy
\texttt{2d-unfolding/HANDOFF\_bkg\_negweight/negweight\_toy.py}, which \emph{is} committed and is
fully deterministic --- literal seeds \texttt{[0, 1, 2]}, \texttt{RandomState(seed)} sampling and
\texttt{LGBMClassifier(random\_state=seed)}, no data and no cluster inputs. It was re-run on
2026-08-21 and reproduced all four exactly; the receipt is
\texttt{docs/orchestration/receipts/RECEIPT-negweight-toy-20260821.json}. They are therefore no
longer struck. The distinction the two classes make is \emph{attestability}, not importance: a toy
needs a re-run, real data needs its products to exist somewhere durable.

\subsection{Estimator}
\label{sec:negweight-estimator}

Step 1 of OmniFold trains a classifier that separates the simulated
reconstruction (class~0, weighted density $S$) from the measured sample
(class~1, weighted density $\rho_1$), then reweights the simulation by the
learned density ratio $r(x)=p(x)/(1-p(x))$, whose population optimum is
$r^{\star}(x)=\rho_1(x)/S(x)$. The background treatment is entirely a choice of
$\rho_1$:
\begin{itemize}
  \item the binned purity weight sets
    $\rho_1(x)=D(x)\,w_{\mathrm{pur}}(x)=\max(0,\,D(x)-B(x))$, realized by a
    per-reco-bin scalar applied to each data event, with $D$ the data density
    and $B$ the POT-scaled background density;
  \item negative-weight injection appends the background-simulation
    reconstruction events to the measured side with weight
    $-w_{\mathrm{bkg}}\,\times\,(\text{POT scale})$, leaving the data events at
    $+1$, so that $\rho_1(x)=D(x)-B(x)$ \emph{continuously}, with no analysis
    binning.
\end{itemize}
Both give the same background-subtracted estimand
$r^{\star}(x)=(D(x)-B(x))/S(x)$; they differ only in whether the subtraction is
performed per reco bin or event by event in the full reconstructed phase space.
The injected background is restricted to the same reconstructed fiducial window
as the data, so the two constructions act on identical support.

\paragraph{The floor.} Where the background locally exceeds the data,
$D-B<0$, the binned method floors the correction at zero through the
$\max(0,\cdot)$. In the unbinned case the classifier output $p\in(0,1)$ forces
$r=p/(1-p)\ge 0$, so a negative reweight cannot arise; the non-negativity
is therefore structural, while the value at the floor follows from the
bounded output and ordinary regularization rather than from a minimizer of
the fit (\S\ref{sec:negweight-binned}). The fit is not well-posed in such a region,
however --- the class-1 weighted mass is locally negative --- so the principled
treatment there is a positive-target prescription. In this analysis we use the
``Stay Positive'' refinement of Ref.~\cite{Nachman:2025staypositive}, followed
by an explicit zero clip of the learned refinement factor, before step-1
training. Whether the plain injection and this clipped refinement coincide in
practice depends on the classifier.
With the histogrammed gradient-boosted density estimator used for \emph{this
comparison} they agree to the median of $\dead{\nwMedianBin}$ reported below, at
the sub-percent genuine-background scale of this measurement. Note that
\texttt{hist} is neither the backend of the production central value
(\texttt{exact}) nor of its covariance ensemble (\texttt{lgbm}); see
\S\ref{sec:method}. This validation therefore establishes the equivalence on a
third backend, and its transfer to the production configuration is an assumption
rather than a demonstration. With an
\emph{exact} (tree-boosted) classifier the raw signed weights instead
destabilize training --- a full low-$p_T$ row of the cross section inflates by
up to a factor ${\sim}\dead{\nwSpRawMax}$ where the local background is large --- and
the clipped Stay-Positive refinement removes the pathology, recovering
agreement with the binned purity to a median of $\dead{\nwSpMedian}$
($\dead{\nwSpWithin}$ bins within \SI{5}{\percent}). This clipped refinement is
therefore what makes the unbinned subtraction robust across estimators, and it
matters most for the larger, more structured backgrounds of the
extended-fiducial analysis.

\subsection{Binned limit: exact reduction to the purity correction}
\label{sec:negweight-binned}

Restrict the step-1 classifier to the sieve of functions piecewise-constant on
the reconstructed analysis bins $\{b\}$ --- the population-level binned limit
that gives the OmniFold/Richardson--Lucy correspondence
of Refs.~\cite{Andreassen:2019cjw,Andreassen:2021zzk}. Within a bin the
classifier output collapses to a single value $p$, and, using the same support
and unnormalized class weights in both constructions, the weighted
cross-entropy of step~1 separates into independent per-bin objectives
\begin{equation}
  \label{eq:negweight-Lb}
  L_b(p) = -W_1\,\ln p - W_0\,\ln(1-p),
\end{equation}
with $W_0=S(b)>0$ the simulated-reconstruction weight of bin $b$ --- the bin
integral of the class-0 density $S$ --- and $W_1$ the signed class-1 weight.
Under negative-weight injection the measured side carries the data at $+1$ and
the background simulation at negative weight, so $W_1=D(b)-B(b)$, writing $D(b)$
and $B(b)$ for the bin masses of $D$ and $B$ (equivalently $N_{\mathrm{data}}(b)$
and $N_{\mathrm{bkg}}(b)$ of \S\ref{sec:negweight-estimator}).

For every bin with $D(b)>B(b)$ the objective~\eqref{eq:negweight-Lb} is
strictly convex on $(0,1)$ and has the unique interior minimizer
$p=W_1/(W_0+W_1)$, so that
\begin{equation}
  \label{eq:negweight-rb}
  r(b) = \frac{p}{1-p} = \frac{W_1}{W_0} = \frac{D(b)-B(b)}{S(b)}.
\end{equation}
This is exactly the reweight produced under the same sieve by the purity
construction, for which the class-1 weight
$W_1=D(b)\,w_{\mathrm{pur}}(b)=\max(0,\,D(b)-B(b))$ coincides with $D(b)-B(b)$ on
this domain. At $D(b)=B(b)$ the infimum lies at the boundary $p\to0^+$ and
$r(b)\to0$, again giving the purity result when that boundary is admitted.
Thus, at the population optimum and under the common binned sieve, the two
background treatments agree bin by bin. Composed with the binned
OmniFold/Richardson--Lucy correspondence, this is precisely the traditional
subtract-background-then-unfold update. Finite, regularized classifiers only
approximate this identity; their residual differences are tested below.

Where $D(b)<B(b)$ the reduction does not extend. With $W_1<0$ the objective
$L_b(p)=|W_1|\,\ln p-W_0\,\ln(1-p)$ tends to $-\infty$ as $p\to0^+$: it is
unbounded below and has no minimizer. A minimizing sequence has $r(b)\to0$,
but adopting $\max(0,\cdot)$ is a prescription, not the value of a minimizer;
the behavior of a finite regularized classifier is estimator-dependent. The
Stay-Positive refinement remains well behaved in this setting but can retain
negative refined weights where the target density is genuinely negative. The
zero clip used here is therefore an additional analysis prescription chosen to
match the physical purity floor. Ref.~\cite{Nachman:2025staypositive} documents
the same divergence of the reweighting loss where the target density turns
negative; signed binary cross-entropy exhibits the analogous pathology in the
sWeights context~\cite{GlazierTyson:2024}.

The three ingredients --- the binned $D-B$ update, negative-weight classification
of the measured side, and the binned OmniFold/Richardson--Lucy correspondence ---
already appear together in Ref.~\cite{Andreassen:2021zzk}. Equation
\eqref{eq:negweight-rb} makes their immediate binwise corollary, and its precise
domain of validity, explicit for this analysis; it is a consistency check rather
than a separate method claim.

\subsection{Closed-form toy}
\label{sec:negweight-toy}

The estimator is verified on a one-dimensional problem with a known answer:
signal simulation drawn from $\mathcal{N}(0,1)$, a data signal from
$\mathcal{N}(0.3,1.1)$ with equal normalization, and a background
$\mathcal{N}(2,0.7)$ contributing \SI{16.7}{\percent} of the pseudo-data, so
that the target reweight $(D-B)/S=\phi(x;0.3,1.1)/\phi(x;0,1)$ is a smooth
positive function. A gradient-boosted classifier trained with the signed-weight
construction recovers this ratio to a median \SI{\nwToyNeg}{\percent} relative
error --- matching the binned purity construction's \SI{\nwToyPur}{\percent} on
the identical problem, and reproducing it bin-for-bin --- while omitting the
subtraction leaves a \SI{\nwToyBias}{\percent} bias from the residual
background. The result is stable across classifier seeds
($\pm\SI{\nwToySeed}{\percent}$), and the negative sample weights introduce no
numerical instability. In a stress configuration with the background moved into
the sparse signal tail and deliberately over-subtracted, the learned reweight
floors cleanly at zero over the region where $B>D$, as anticipated above.

\subsection{Comparison on the 2D measurement}
\label{sec:negweight-2d}

Table~\ref{tab:negweight} compares the two background treatments on the
ME~FHC data, unfolded with identical settings and a common classifier
seed so that the difference isolates the background treatment. The two agree to
\dead{\SI{\nwPctTot}{\percent}} on the total cross section --- an order of magnitude
inside the ${\sim}\SI{3}{\percent}$ purity correction itself --- with a median
per-bin ratio of $\dead{\nwMedianBin}$ and a \dead{\SI{\nwRmsBin}{\percent}} RMS scatter
across the plane (Fig.~\ref{fig:negweight-ratio}).
The only bin differing by more than a few percent is a single low-$p_T$,
high-background edge bin (\dead{\SI{\nwWorstBin}{\percent}}), exactly where the
background concentrates and the subtraction is largest. The unbinned
subtraction thus reproduces the adopted binned correction wherever the
correction is well determined, and departs from it only in the regime where the
binned $\max(0,\cdot)$ floor and the structural floor of
\S\ref{sec:negweight-estimator} are themselves the operative choice.

\begin{table}[t]
  \centering
  \caption{Binned purity versus unbinned negative-weight background
    subtraction on the 2D ME~FHC measurement (common estimator and seed).
    The per-bin figures are the negative-weight/purity ratio over the
    populated cross-section bins.}
  \label{tab:negweight}
  \begin{tabular}{lcc}
    \hline\hline
    & binned purity & negative-weight \\
    \hline
    total $\sigma$ (cm$^2$/nucleon) & $\dead{\nwSigPur}$ & $\dead{\nwSigNeg}$ \\
    total ratio (neg/pur)           & \multicolumn{2}{c}{$\dead{\nwRatioTot}$ (\dead{\SI{\nwPctTot}{\percent}})} \\
    median per-bin ratio            & \multicolumn{2}{c}{$\dead{\nwMedianBin}$} \\
    per-bin RMS deviation           & \multicolumn{2}{c}{\dead{\SI{\nwRmsBin}{\percent}}} \\
    largest single-bin deviation    & \multicolumn{2}{c}{\dead{\SI{\nwWorstBin}{\percent}} (low $p_T$, high background)} \\
    \hline\hline
  \end{tabular}
\end{table}

\begin{figure}[t]
  \centering
  \includegraphics[width=0.62\textwidth]{negweight_ratio_2d}
  \caption{Per-bin ratio of the negative-weight to the binned-purity
    cross section on the 2D ME~FHC measurement (common estimator and seed),
    over the populated bins. The two agree to a median of $\dead{\nwMedianBin}$ across
    the plane (near-white); the only departure beyond a few percent is the
    single low-$p_T$, high-background edge bin, where the background is largest
    and the $\max(0,\cdot)$/structural floor of \S\ref{sec:negweight-estimator}
    is the operative choice.}
  \label{fig:negweight-ratio}
\end{figure}

\subsection{Propagation through systematics and statistics}
\label{sec:negweight-uq}

A central-value match is necessary but not sufficient: to replace the binned
correction the unbinned subtraction must also reproduce the \emph{uncertainty}
budget. The full 2D ME~FHC covariance was therefore rebuilt end to end under
negative-weight injection --- the systematic ensemble of 187 universes, with the
genuine background varied consistently within each universe rather than frozen
at nominal, and a statistical bootstrap --- and compared to the binned-purity
construction with identical settings. The two agree at the level the $(D-B)/S$
identity predicts: the systematic covariance matches in overall magnitude to a
$\sqrt{\mathrm{tr}}$ ratio of $\dead{\nwSystRatio}$ (negative-weight/purity), and the
statistical covariance, compared on matched bootstrap seeds, to $\dead{\nwStatRatio}$.
The residual differences --- \dead{\SI{\nwSystResid}{\percent}} and
\dead{\SI{\nwStatResid}{\percent}} --- \textbf{struck with their parents: these are DERIVED from
$\dead{\nwSystRatio}$ and $\dead{\nwStatRatio}$ and are now DEFINED FROM THEM MECHANICALLY, so
they cannot survive their parents' retraction again} --- consistent with
the ``to within \SI{2}{\percent}'' stated in \S\ref{sec:method} --- are bootstrap and
classifier-seed sampling noise, not a method bias. Negative-weight injection thus
reproduces the central cross section \emph{per bin} (median $1.000$, RMS
\SI{1.4}{\percent}, worst bin $-\SI{12.6}{\percent}$) and its covariance \emph{in
aggregate}. The covariance comparison is two trace ratios; per-bin, correlation and
eigenmode equivalence are not established.

The equivalence at the 2D scale --- in the central value, the closed-form
estimator recovery, and now the full systematic and statistical covariance ---
together with the Stay-Positive refinement that keeps the estimator robust on
the larger backgrounds of the extended-fiducial analysis, establishes
negative-weight injection as a fully validated unbinned replacement for the
binned purity correction. The binned purity is retained as the default for the
reproduction measurement of this note, with negative-weight injection carried as
the validated unbinned cross-check and as the \emph{designated} baseline
subtraction for the extended-fiducial analysis (\S\ref{sec:fps}), where its
binning-independence and continuous floor are essential.

\paragraph{What the quoted extended-fiducial number actually used.}
\textbf{The $\sigma_{\mathrm{ext}}=\SI{\fpsSigma}{cm^2/nucleon}$ reported in
\S\ref{sec:fps} was not produced with negative-weight injection}, and an earlier
version of this appendix implied that it was. The product is dated: it comes from
the FPS MEFHC battery of \textbf{2026-06-10}, job \texttt{54244120}
(\texttt{VALIDATION\_LEDGER.md}, ``FPS MEFHC battery''), whereas the
negative-weight driver code did not land until \textbf{2026-07-11}
(\texttt{cf8a4a6}) and the \texttt{-{}-bkg-mode negweight-refined} launcher flag
not until 2026-07-20. Both FPS unfold launchers omitted \texttt{-{}-bkg-mode}
entirely, so the ten endpoint unfolds ran on the driver default
\texttt{-{}-bkg-mode=purity}; they are recorded as \textbf{purity controls, not
publication inputs}
(\texttt{nd-unfolding/uq\_fps/corrected/FPS\_UQ\_CORRECTED\_STATE.md}). The
negweight-refined FPS production has not been run.

The \emph{value} is unaffected --- \SI{\fpsSigma}{cm^2/nucleon} remains the
ledger-verified central number from that gated battery, and this correction does
not change it. What is withdrawn is the \emph{footing}: negative-weight injection
is the intended baseline for the extended-fiducial analysis, not the one behind
the number quoted here. (Audit finding J22,
\texttt{docs/orchestration/AUDIT-FINDINGS-20260731.md}, tier A.)

\subsection{Which footing each production chain stands on}
\label{sec:negweight-footing}

The two constructions above are validated against one another; that does not make
the choice between them arbitrary, and the two production chains of this analysis
do not make it the same way. This subsection records which chain stands on which
footing, on what measured basis, and --- because over-reading that basis is the
reason the subsection is written at all --- what it does \emph{not} establish.

\paragraph{The standard \(N\)-dimensional chain is purity-footed.}
The whole standard (five-dimensional) chain --- the central unfold, the $169$
vertical-band unfolds, the $18$ detector unfolds and the $10$ active-lateral
endpoint unfolds --- uses the binned purity correction of
\S\ref{sec:method-algo}. It is positively identified as purity-footed rather
than merely unstamped, and the evidence differs between the unfolds that already
exist and the chain as it is run today. \emph{For the existing endpoint unfolds},
produced in 2026-07, the identification is by elimination: neither launcher that
could have produced them passes \texttt{-{}-bkg-mode}, the driver's default for
that flag is \texttt{purity}, and the driver announces its mode only on the
negative-weight branches --- so the lateral logs carry the purity branch's own
signature line and no mode announcement anywhere. Both the flag and those
announcements predate the files, so the absence of an announcement is
informative rather than a version gap. \emph{For the chain as it now stands},
the footing is no longer inferred at all: it is pinned in code
(\texttt{nd-unfolding/p4\_lib.py}, \texttt{STANDARD\_BKG\_MODE}, whose
configuration object refuses to validate against any other mode), the canonical
launcher passes \texttt{-{}-bkg-mode} explicitly from that pinned configuration
rather than inheriting the default, and the value is stamped into every endpoint
receipt as it is written. That change is a provenance change and a physics
no-op, since the value it asserts is the default it replaces; its effect is that
a future component's footing is provable from the component instead of from the
absence of a flag.

\textbf{This is a recorded choice, not a silent default.} It was raised as a
question --- whether the locked instruction that $N$-D production runs
\texttt{negweight-refined} reaches the standard-phase-space chain, or only its
extended-fiducial form --- and answered in favour of purity on 2026-08-07 and
again, after the alternative was costed, on 2026-08-15
(\texttt{docs/orchestration/RUNBOOK-20260807-gbdt-closeout.md} \S2.1--2.2;
\texttt{docs/orchestration/DECISION-20260815-joseph-oi6-oi8-oi126.md}). Two
grounds carried it. First, mixing footings inside one covariance is worse than
either consistent choice: dropping a negative-weight-footed lateral block into a
purity-footed central and vertical set is a footing mismatch, not an upgrade.
Second, the alternative required re-running every unfold named above, which would
have invalidated the flux-corrected (J28) covariance adopted immediately
beforehand --- and it is that consequence, not the cost, that decided it. The
decision carries an explicit obligation to revisit the footing before
submission. It is recorded as such, and nothing in this subsection discharges it.

\paragraph{So the published pair is a measurement and a matched control at
different footings.} The extended-fiducial lane's designated publication footing
is negative-weight injection with the Stay-Positive refinement
(\texttt{nd-unfolding/fps\_provenance.py}, \texttt{PUBLICATION\_BKG\_MODE}),
enforced there by a gate that fails closed on any endpoint whose recorded mode
differs or is missing; what the number quoted in \S\ref{sec:fps} actually used is
the separate matter stated immediately above. The standard chain is purity-footed
throughout. Each side is therefore internally consistent in its footing, which is
what we claim, rather than uniform across both, which we do not.

\paragraph{The measured size of the choice.}
It has been measured at full statistics with the adopted estimator on the 2D
ME~FHC measurement, and it is one to two percent on the covariance. The
$187$-universe systematic covariance agrees to a $\sqrt{\mathrm{Tr}}$ ratio of
$\dead{\nwSystRatio}$, from \dead{\SI{2.9828e-39}{cm^2/nucleon}} negative-weight against
\dead{\SI{3.0242e-39}{cm^2/nucleon}} purity; the statistical covariance on matched
bootstrap seeds agrees to $\dead{\nwStatRatio}$, from \dead{\SI{1.7260e-40}{cm^2/nucleon}}
against \dead{\SI{1.7576e-40}{cm^2/nucleon}}; and the real-data totals agree to
\dead{\SI{\nwPctTot}{\percent}}, with per-bin median $\dead{\nwMedianBin}$ and
\dead{\SI{\nwRmsBin}{\percent}} RMS. These are the comparisons of
\S\ref{sec:negweight-2d} and \S\ref{sec:negweight-uq}, restated here with the
operands of each ratio so that the ratios can be checked rather than taken.

\paragraph{Why the agreement is structural rather than two lucky universes.}
\S\ref{sec:negweight-binned} makes the two constructions the same estimand,
$\rho_1=D-B$, so what separates them is the finite-classifier residual, and at
this measurement's sub-percent genuine-background scale that residual is a shift
of order \SI{0.1}{\percent} per universe. A systematic covariance is the spread
\emph{across} universes about their own mean. A footing change that displaces
every universe by the same small amount therefore leaves that spread almost
unchanged: the two covariances agree by construction, not by coincidence, which
is why the $187$-universe ratio lands at $\dead{\nwSystRatio}$ rather than anywhere near
the size of the purity correction itself (${\sim}\SI{3}{\percent}$). A reader is
entitled to ask how a comparison on a handful of universes can support a
statement about a $187$-universe object; the displacement argument is the reason
it can, and the full $187$-universe rebuild of \S\ref{sec:negweight-uq} is that
reasoning's test at 2D, which it passes.

\paragraph{What this does not establish, and must not be edited into
establishing.} \textbf{There is no full five-dimensional $187$-universe
both-footing comparison at the publication configuration} ($5$ iterations,
\texttt{lgbm}). The five-dimensional evidence is a two-universe spot check at one
iteration on the \texttt{hist} backend --- on the on-disk five-dimensional
universe file, \texttt{2p2h:0} gives $3.014\times10^{-38}$ under both footings
with per-bin median $1.0005$, and \texttt{MaCCQE:0} differs by
\SI{-0.03}{\percent} with median $0.9990$ --- together with the displacement
argument above and the full-statistics 2D result. That is ample for quoting the
standard chain on purity and saying so, which is what this subsection does. It is
\emph{not} ample for the stronger statement that the background footing is
irrelevant in five dimensions, and \textbf{that statement should not be
introduced here by a later edit}: it would require the comparison named at the
head of this paragraph, and that comparison has not been run. The same limit
applies to the one estimator caveat: raw negative weights destabilise the
\texttt{exact} backend (\S\ref{sec:negweight-estimator}) and the standard chain
runs \texttt{lgbm}, so neither footing risks that pathology here --- which is a
reason the choice is safe, not evidence that it is immaterial.

\paragraph{The footing question is settled; the adoption gate is a different
one.} \textbf{Nothing here adopts the standard five-dimensional covariance}, and
this subsection must not be read as doing so. That covariance remains a candidate
(\S\ref{sec:uq-combine}) because its selection-complete active-lateral component
is not built: the ten purity-footed endpoint unfolds named above exist but carry
no completion receipt and were produced by a since-retired launcher, and building
the component from them is gated by an internal verifier that has not authorised
it. That verifier's current verdict
(\texttt{docs/orchestration/runs/standard-p4-verifier/}\allowbreak
\texttt{20260810T012645Z-repair7-verdict.json}) is \texttt{BLOCK}, with fourteen
defects outstanding, and it does not authorise the covariance-construction
stages. The two questions are independent, and settling the footing removes a
decision that was blocking while leaving every other precondition on that
component exactly where it stood.
